\chapter{Complex Numbers}\label{ch:g12:complex}

The \omterm{def:g10:algebra:equation}{equation} $x^2 = -1$ has no real solution. Enlarging $\R$ with a single
new number $\iu$ whose square is $-1$ produces the field $\C$ of \omterm{def:g12:complex:def}{complex
numbers}, in which \emph{every} polynomial \omterm{def:g10:algebra:equation}{equation} has solutions --- and
whose arithmetic encodes the geometry of the plane: translations, rotations
and scalings become additions and multiplications.

\section{Algebraic form}

\begin{definition}[Complex numbers]\label{def:g12:complex:def}
The set of \emph{complex numbers}\index{complex number} is
\[
\C = \{a + \iu b : a, b \in \R\},
\]
where $\iu$ is a symbol subject to the single rule $\iu^2 = -1$; addition and
multiplication are performed as for real polynomials in $\iu$, reducing
$\iu^2$ to $-1$. The reals $a$ and $b$ are the \emph{real part}\index{real part}
$\Rea(z)$ and \emph{imaginary part}\index{imaginary part} $\Ima(z)$ of $z = a + \iu b$; this
expression is the \emph{algebraic form}\index{algebraic form} of $z$, and it is unique:
$a + \iu b = a' + \iu b'$ if and only if $a = a'$ and $b = b'$.
\end{definition}

\begin{example}
$(2 + 3\iu)(1 - \iu) = 2 - 2\iu + 3\iu - 3\iu^2 = 5 + \iu$.
\end{example}

\begin{definition}[Conjugate]\label{def:g12:complex:conjugate}
The \emph{conjugate}\index{conjugate} of $z = a + \iu b$ is
$\conj{z} = a - \iu b$.
\end{definition}

\begin{proposition}[Properties of conjugation]\label{prop:g12:complex:conjugate}
For all $z, w \in \C$:
\[
\conj{z + w} = \conj z + \conj w, \quad
\conj{zw} = \conj z\,\conj w, \quad
\conj{\conj z} = z, \quad
z + \conj z = 2\Rea(z), \quad
z - \conj z = 2\iu\,\Ima(z),
\]
and $z\conj z = a^2 + b^2 \geq 0$ for $z = a + \iu b$. Moreover $z \in \R$
if and only if $z = \conj z$.
\end{proposition}

\begin{proof}
All are direct computations from the \omterm{def:g12:complex:def}{algebraic forms}; for instance
$z\conj z = (a + \iu b)(a - \iu b) = a^2 - (\iu b)^2 = a^2 + b^2$.
\end{proof}

\begin{method}[Dividing complex numbers]\label{met:g12:complex:division}
To write a quotient in \omterm{def:g12:complex:def}{algebraic form}, multiply numerator and denominator by
the \omterm{def:g12:complex:conjugate}{conjugate} of the denominator:
\[
\frac{1}{2 + \iu} = \frac{2 - \iu}{(2+\iu)(2-\iu)} = \frac{2 - \iu}{5}
= \frac25 - \frac15\iu .
\]
Every nonzero \omterm{def:g12:complex:def}{complex number} therefore has an inverse: $\C$ is a
\emph{field}.
\end{method}

\section{The complex plane, modulus and argument}

\begin{definition}[Affix, modulus]\label{def:g12:complex:modulus}
In the plane with \omterm{def:g10:coordgeom:system}{orthonormal} \omterm{def:g10:coordgeom:system}{coordinates}, the point $M(a, b)$ is the
\emph{\omterm{def:g10:functions:function}{image}} of $z = a + \iu b$, and $z$ is the \emph{affix}\index{affix} of $M$. The
\emph{modulus}\index{modulus} of $z$ is the distance from $M$ to the origin:
\[
\abs{z} = \sqrt{a^2 + b^2} = \sqrt{z\conj z}.
\]
\end{definition}

\begin{omfigure}
\begin{tikzpicture}[scale=1.5]
  \draw[->] (-0.7,0) -- (2.9,0) node[right] {$x$};
  \draw[->] (0,-1.9) -- (0,2.1) node[above] {$y$};
  \draw[dashed] (2,1.5) -- (2,0) node[below right] {$a$};
  \draw[dashed] (2,1.5) -- (0,1.5) node[left] {$b$};
  \draw[dashed] (2,-1.5) -- (2,0);
  \draw[-Stealth,omDef,thick] (0,0) -- (2,1.5)
    node[midway, above left, font=\small] {$\abs z$};
  \draw[-Stealth,omThm,thick] (0,0) -- (2,-1.5);
  \fill (2,1.5) circle (1pt) node[above right] {$M(z = a + \iu b)$};
  \fill (2,-1.5) circle (1pt) node[below right] {$\conj z = a - \iu b$};
\end{tikzpicture}

{\small The point $M(a,b)$ of \omterm{def:g12:complex:modulus}{affix} $z$: the \omterm{def:g12:complex:modulus}{modulus} $\abs z$ is the
length $OM$, and the \omterm{def:g12:complex:conjugate}{conjugate} $\conj z$ is the reflection of $z$ in the
real axis.}
\end{omfigure}

\begin{proposition}[Properties of the modulus]\label{prop:g12:complex:modulus}
For $z, w \in \C$:
\[
\abs{zw} = \abs z\,\abs w, \qquad
\abs{\tfrac zw} = \tfrac{\abs z}{\abs w}\ (w \neq 0), \qquad
\abs{\conj z} = \abs z, \qquad
\abs{z + w} \leq \abs{z} + \abs{w} \ \text{(triangle inequality)}.
\]
Moreover $\abs{z_B - z_A}$ is the distance between the points of \omterm{def:g12:complex:modulus}{affixes}
$z_A$ and $z_B$.
\end{proposition}

\begin{proof}
$\abs{zw}^2 = zw\,\conj{zw} = z\conj z\, w \conj w = \abs z^2 \abs w^2$
gives the first; the second and third are similar. For the triangle
inequality:
\[
\abs{z+w}^2 = (z+w)(\conj z + \conj w)
= \abs z^2 + \abs w^2 + 2\Rea(z \conj w)
\leq \abs z^2 + \abs w^2 + 2\abs{z\conj w}
= \bigl(\abs z + \abs w\bigr)^2,
\]
using $\Rea(u) \leq \abs u$. The distance statement is the Pythagorean
formula applied to the \omterm{def:g10:coordgeom:system}{coordinates} of $z_B - z_A$.
\end{proof}

\begin{definition}[Argument, trigonometric and exponential forms]\label{def:g12:complex:argument}
Let $z \neq 0$. An \emph{argument}\index{argument} of $z$, written
$\Arg(z)$, is any measure $\theta$ of the angle from the positive real axis
to the ray $OM$; it is defined up to $2k\pi$. Writing $r = \abs z$,
\[
z = r(\cos\theta + \iu\sin\theta)
\qquad\text{(trigonometric form)}.
\]
Introducing the notation\index{exponential form}
\[
\eu^{\iu\theta} = \cos\theta + \iu\sin\theta,
\]
this becomes the \emph{exponential form} $z = r\,\eu^{\iu\theta}$.
\end{definition}

\begin{omfigure}
\begin{tikzpicture}[scale=1.5]
  \draw[->] (-0.7,0) -- (2.4,0) node[right] {$x$};
  \draw[->] (0,-0.5) -- (0,1.9) node[above] {$y$};
  \draw[gray, dashed] (1.9,0) arc (0:75:1.9);
  \draw[-Stealth,omDef,thick] (0,0) -- (40:1.9)
    node[pos=0.62, below right=1pt, font=\small] {$r = \abs z$};
  \fill (40:1.9) circle (1pt) node[above right] {$z = r\,\eu^{\iu\theta}$};
  \draw[->,omThm,thick] (0.45,0) arc (0:40:0.45);
  \node[omThm] at (20:0.65) {$\theta$};
\end{tikzpicture}

{\small \omterm{def:g12:complex:argument}{Exponential form}: $z$ is determined by its distance $r$ to the
origin and the angle $\theta$ from the positive real axis.}
\end{omfigure}

\begin{theorem}\label{thm:g12:complex:expform}
For all $\theta, \theta' \in \R$:
\[
\eu^{\iu\theta}\,\eu^{\iu\theta'} = \eu^{\iu(\theta + \theta')},
\qquad
\conj{\eu^{\iu\theta}} = \eu^{-\iu\theta},
\qquad
\abs{\eu^{\iu\theta}} = 1 .
\]
Consequently, for $z = r\eu^{\iu\theta}$ and $z' = r'\eu^{\iu\theta'}$
nonzero:
\[
zz' = rr'\,\eu^{\iu(\theta+\theta')},
\qquad
\frac{z}{z'} = \frac{r}{r'}\,\eu^{\iu(\theta - \theta')} :
\]
\emph{moduli multiply, \omterm{def:g12:complex:argument}{arguments} add.}
\end{theorem}

\begin{proof}
The first identity \emph{is} the pair of addition formulas
(\cref{prop:g12:trigo:addition}):
\[
(\cos\theta + \iu\sin\theta)(\cos\theta' + \iu\sin\theta')
= (\cos\theta\cos\theta' - \sin\theta\sin\theta')
+ \iu(\sin\theta\cos\theta' + \cos\theta\sin\theta').
\]
The rest follows by direct computation.
\end{proof}

\begin{corollary}[De Moivre's formula]\label{cor:g12:complex:demoivre}
\index{De Moivre's formula}
For $\theta \in \R$ and $n \in \Z$:
$\bigl(\cos\theta + \iu\sin\theta\bigr)^n = \cos n\theta + \iu \sin n\theta$.
\end{corollary}

\begin{proof}
Induction on $n \geq 0$ using \cref{thm:g12:complex:expform}; extend to
$n < 0$ by taking inverses.
\end{proof}

\begin{method}[Switching between forms]\label{met:g12:complex:forms}
From algebraic to exponential: compute $r = \abs z$, then find $\theta$ with
$\cos\theta = \frac ar$, $\sin\theta = \frac br$ (locate the correct quadrant
before invoking an inverse \omterm{def:g12:trigo:cossin}{cosine}). From exponential to algebraic: expand
$r\cos\theta + \iu\, r\sin\theta$. Use the \omterm{def:g12:complex:def}{algebraic form} for sums, the
\omterm{def:g12:complex:argument}{exponential form} for products, powers and quotients.
\end{method}

\begin{example}
$z = 1 + \iu$: $\abs z = \sqrt2$ and $\cos\theta = \sin\theta =
\frac{1}{\sqrt2}$, so $\theta = \frac\pi4$ and $z = \sqrt2\,\eu^{\iu\pi/4}$.
Hence $z^{20} = 2^{10}\,\eu^{5\iu\pi} = -1024$.
\end{example}

\section{Second-degree equations and roots of unity}

\begin{theorem}[Quadratic equations with real coefficients]\label{thm:g12:complex:quadratic}
Let $a, b, c \in \R$, $a \neq 0$, and $\Delta = b^2 - 4ac$. The \omterm{def:g10:algebra:equation}{equation}
$az^2 + bz + c = 0$ has solutions in $\C$:
\[
\begin{aligned}
&\Delta > 0:\ z = \frac{-b \pm \sqrt{\Delta}}{2a} \ (\text{two real}); \qquad
\Delta = 0:\ z = \frac{-b}{2a} \ (\text{double}); \\
&\Delta < 0:\ z = \frac{-b \pm \iu\sqrt{-\Delta}}{2a} \ (\text{two conjugate}).
\end{aligned}
\]
\end{theorem}

\begin{proof}
Complete the square:
$az^2 + bz + c = a\left(\left(z + \frac{b}{2a}\right)^2 -
\frac{\Delta}{4a^2}\right)$. If $\Delta < 0$, then
$\frac{\Delta}{4a^2} = \left(\iu\frac{\sqrt{-\Delta}}{2a}\right)^2$, and the
difference of squares factors as usual.
\end{proof}

\begin{definition}[Roots of unity]\label{def:g12:complex:rootsofunity}
For $n \geq 1$, the \emph{$n$-th roots of unity}\index{roots of unity} are
the solutions of $z^n = 1$.
\end{definition}

\begin{theorem}\label{thm:g12:complex:rootsofunity}
The \omterm{def:g10:algebra:equation}{equation} $z^n = 1$ has exactly $n$ solutions:
\[
\omega_k = \eu^{2\iu k\pi/n}, \qquad k = 0, 1, \dots, n-1 .
\]
Their \omterm{def:g10:functions:function}{images} are the vertices of a regular $n$-gon inscribed in the unit
circle, and for $n \geq 2$ their sum is $0$.
\end{theorem}

\begin{proof}
Each $\omega_k$ satisfies $\omega_k^n = \eu^{2\iu k\pi} = 1$, and they are
pairwise distinct since their \omterm{def:g12:complex:argument}{arguments} lie in $\intco{0}{2\pi}$.
Conversely, if $z^n = 1$ then $\abs z^n = 1$, so $\abs z = 1$ (positive
real) and $z = \eu^{\iu\theta}$ with $n\theta \equiv 0 \pmod{2\pi}$,
\emph{i.e.} $\theta = \frac{2k\pi}{n}$; reducing $k$ modulo $n$ lands among
the $\omega_k$. The vertices are equally spaced by angle $\frac{2\pi}{n}$:
a regular $n$-gon. Finally, with $\omega = \omega_1$,
$\omega \neq 1$ and the geometric sum gives
\[
\sum_{k=0}^{n-1}\omega_k = \sum_{k=0}^{n-1}\omega^k
= \frac{\omega^n - 1}{\omega - 1} = 0 . \qedhere
\]
\end{proof}

\begin{omfigure}
\begin{tikzpicture}[scale=1.7]
  \draw[->] (-1.3,0) -- (1.45,0) node[right] {$x$};
  \draw[->] (0,-1.3) -- (0,1.4) node[above] {$y$};
  \draw[gray] (0,0) circle (1);
  \draw[omDef] (0:1) -- (72:1) -- (144:1) -- (216:1) -- (288:1) -- cycle;
  \draw[->,omThm,thick] (0.3,0) arc (0:72:0.3);
  \node[omThm, font=\small] at (36:0.52) {$\frac{2\pi}{5}$};
  \fill (0:1) circle (1pt) node[below right] {$1$};
  \fill (72:1) circle (1pt) node[above right] {$\omega$};
  \fill (144:1) circle (1pt) node[above left] {$\omega^2$};
  \fill (216:1) circle (1pt) node[below left] {$\omega^3$};
  \fill (288:1) circle (1pt) node[below right] {$\omega^4$};
\end{tikzpicture}

{\small The fifth \omterm{def:g12:complex:rootsofunity}{roots of unity}, $\omega = \eu^{2\iu\pi/5}$: the
vertices of a regular pentagon inscribed in the unit circle
(\cref{exo:g12:complex:10}).}
\end{omfigure}

\begin{example}
The cube roots of unity are $1$, $j = \eu^{2\iu\pi/3} =
-\frac12 + \iu\frac{\sqrt3}{2}$ and $\conj j = j^2$, satisfying
$1 + j + j^2 = 0$.
\end{example}

\section{Complex numbers and plane geometry}

\begin{proposition}[Geometric dictionary]\label{prop:g12:complex:geometry}
Let $A, B, C, D$ be points of \omterm{def:g12:complex:modulus}{affixes} $a, b, c, d$ ($a \neq b$,
$c \neq d$).
\begin{enumerate}
  \item The translation by the \omterm{def:g11:vect:vector}{vector} of \omterm{def:g12:complex:modulus}{affix} $t$ is $z \mapsto z + t$.
  \item The rotation of center $\omega$ (\omterm{def:g12:complex:modulus}{affix}) and angle $\theta$ is
        $z \mapsto \omega + \eu^{\iu\theta}(z - \omega)$.
  \item The scaling (homothety) of center $\omega$ and ratio
        $k \in \R^*$ is $z \mapsto \omega + k(z - \omega)$.
  \item $\dfrac{d - c}{b - a}$ has \omterm{def:g12:complex:modulus}{modulus} $\dfrac{CD}{AB}$ and \omterm{def:g12:complex:argument}{argument} the
        angle between the \omterm{def:g11:vect:vector}{vectors} $\vect{AB}$ and $\vect{CD}$. In
        particular, $AB \perp CD$ if and only if $\frac{d-c}{b-a}$ is purely
        imaginary, and $A, B, C$ are \omterm{def:g11:vect:collinear}{collinear} ($C \ne A$, $C\ne B$) if and
        only if $\frac{c-a}{b-a} \in \R$.
\end{enumerate}
\end{proposition}

\begin{proof}
(1) and (3) restate the coordinate formulas. For (2), the map
$w \mapsto \eu^{\iu\theta} w$ multiplies moduli by $1$ and adds $\theta$ to
\omterm{def:g12:complex:argument}{arguments}: it is the rotation of angle $\theta$ about the origin; conjugating
by the translation sending $\omega$ to $0$ gives the general center. For
(4), moduli and \omterm{def:g12:complex:argument}{arguments} of quotients subtract
(\cref{thm:g12:complex:expform}).
\end{proof}

\section{Exercises}

\begin{exercise}[$\star$]\label{exo:g12:complex:1}
Write in \omterm{def:g12:complex:def}{algebraic form}:
$(3 - 2\iu)(1 + 4\iu)$, \quad $\dfrac{2 + \iu}{1 - \iu}$, \quad
$\iu^{2027}$.
\end{exercise}

\begin{exercise}[$\star$]\label{exo:g12:complex:2}
Solve in $\C$: (a) $z^2 - 4z + 13 = 0$; (b) $z^2 + z + 1 = 0$. Check in each
case that the two solutions are \omterm{def:g12:complex:conjugate}{conjugates} and compute their product.
\end{exercise}

\begin{exercise}[$\star$]\label{exo:g12:complex:3}
Put in \omterm{def:g12:complex:argument}{exponential form}: $z_1 = -1 + \iu$, $z_2 = \sqrt3 - \iu$, and compute
$\dfrac{z_1}{z_2}$ in exponential then \omterm{def:g12:complex:def}{algebraic form}. Deduce the exact
value of $\cos\dfrac{11\pi}{12}$.
\end{exercise}

\begin{exercise}[$\star\star$]\label{exo:g12:complex:4}
Describe geometrically the set of points $M$ of \omterm{def:g12:complex:modulus}{affix} $z$ such that:
(a) $\abs{z - 2} = \abs{z + \iu}$;\quad
(b) $\abs{z - 1 - \iu} = 3$;\quad
(c) $\dfrac{z - 1}{z + 1}$ is purely imaginary.
\end{exercise}

\begin{exercise}[$\star\star$]\label{exo:g12:complex:5}
Using De Moivre's formula and the binomial theorem, express $\cos 3\theta$
as a polynomial in $\cos\theta$, and $\sin 3\theta$ in terms of
$\sin\theta$.
\end{exercise}

\begin{exercise}[$\star\star$]\label{exo:g12:complex:6}
\emph{(Linearization.)} Using
$\cos\theta = \frac{\eu^{\iu\theta} + \eu^{-\iu\theta}}{2}$, linearize
$\cos^3\theta$ (write it as a combination of $\cos k\theta$), and deduce
$\displaystyle\int_0^{\pi/2}\cos^3\theta\,\dd\theta$.
\end{exercise}

\begin{exercise}[$\star\star$]\label{exo:g12:complex:7}
Let $A$ and $B$ have \omterm{def:g12:complex:modulus}{affixes} $a = 1 + \iu$ and $b = 3 + 2\iu$. Determine the
two points $C$ making the triangle $ABC$ equilateral, as the \omterm{def:g10:functions:function}{images} of $B$
under the rotations of center $A$ and angles $\pm\frac{\pi}{3}$.
\end{exercise}

\begin{exercise}[$\star\star\star$]\label{exo:g12:complex:8}
Solve $z^4 = -4$ in $\C$ and plot the solutions. Factor
$z^4 + 4$ into two quadratic polynomials with real coefficients.
\end{exercise}

\begin{exercise}[$\star\star\star$]\label{exo:g12:complex:9}
For $\theta \in \intoo{0}{2\pi}$ and $n \in \N$, compute
\[
S = 1 + \cos\theta + \cos 2\theta + \dots + \cos n\theta
\]
by summing the geometric progression
$\sum_{k=0}^n \eu^{\iu k\theta}$ and taking \omterm{def:g12:complex:def}{real parts}. (Answer:
$S = \dfrac{\sin\frac{(n+1)\theta}{2}}{\sin\frac{\theta}{2}}
\cos\frac{n\theta}{2}$.)
\end{exercise}

\begin{exercise}[$\star\star\star$]\label{exo:g12:complex:10}
Let $\omega = \eu^{2\iu\pi/5}$.
\begin{enumerate}
  \item Justify that $1 + \omega + \omega^2 + \omega^3 + \omega^4 = 0$.
  \item Let $u = \omega + \omega^4$ and $v = \omega^2 + \omega^3$. Show that
        $u + v = -1$ and $uv = -1$, and deduce that $u = \frac{-1+\sqrt5}{2}$.
  \item Conclude that $\cos\dfrac{2\pi}{5} = \dfrac{\sqrt5 - 1}{4}$.
\end{enumerate}
\end{exercise}

\section{Problem: The pentagon in the mirror of unity}

\begin{problem}[{Weekend problem --- the fifth roots of unity
compute $\cos 72^\circ$ exactly, the golden ratio signs the
result, and multiplication turns out to be rotation}]
\label{pb:g12:complex:1}
Euclid could construct the regular pentagon but the reason
\emph{why} it yields to ruler and compass --- while the humble
$20^\circ$ angle does not --- stayed hidden for two thousand
years. It hides in this chapter: the five fifth \omterm{def:g12:complex:rootsofunity}{roots of unity}
(\cref{thm:g12:complex:rootsofunity}) sum to zero, and that one
line of algebra squeezes $\cos 72^\circ$ out of a quadratic
\omterm{def:g10:algebra:equation}{equation} --- square roots only, hence constructible, and signed,
of course, by the golden ratio. Around this jewel: fluency, de
Moivre's shortcut to trigonometry, and multiplication unmasked as
geometry.

\medskip
\noindent\textbf{Part I --- Fluency.}
\begin{enumerate}
  \item Compute $(2 + \iu)(3 - \iu)$, \
        $\dfrac{1 + \iu}{1 - \iu}$
        (\cref{met:g12:complex:division}), and
        $\abs{3 + 4\iu}$.
  \item Put in \omterm{def:g12:complex:argument}{exponential form}: $1 + \iu$; \ $-2$; \
        $1 - \iu\sqrt3$.
  \item Compute $(1 + \iu)^8$ using the \omterm{def:g12:complex:argument}{exponential form}.
  \item Solve $z^2 = \iu$, then $z^2 - 2z + 5 = 0$
        (\cref{thm:g12:complex:quadratic}).
  \item Geometric dictionary
        (\cref{prop:g12:complex:geometry}): describe the map
        $z \mapsto \iu z$, and the set
        $\abs{z - (1 + \iu)} = 2$.
\end{enumerate}

\medskip
\noindent\textbf{Part II --- Five \omterm{def:g11:quad:discriminant}{roots}, one pentagon.}
Let $\omega = \eu^{2\iu\pi/5}$ and consider
$1, \omega, \omega^2, \omega^3, \omega^4$.
\begin{enumerate}[resume]
  \item Check these are exactly the solutions of $z^5 = 1$,
        and describe the figure they draw in the plane.
  \item Prove that their sum is zero. (Geometric progression
        --- or multiply the sum by $1 - \omega$.)
  \item Take \omterm{def:g12:complex:def}{real parts} and deduce
        \[
        1 + 2\cos\frac{2\pi}{5} + 2\cos\frac{4\pi}{5} = 0 .
        \]
  \item Set $c = \cos\frac{2\pi}{5}$ and use the double-angle
        formula for $\cos\frac{4\pi}{5}$ to convert question 8
        into a quadratic \omterm{def:g10:algebra:equation}{equation} for $c$; solve it and
        conclude
        \[
        \cos 72^\circ = \frac{\sqrt5 - 1}{4} .
        \]
  \item Verify numerically, then let the golden ratio sign its
        work: show $c = \frac{1}{2\varphi}$ where
        $\varphi = \frac{1 + \sqrt5}{2}$
        (\cref{pb:g10:algebra:1}) --- and state the famous
        geometric echo: in a regular pentagon, the diagonal is
        $\varphi$ times the side.
\end{enumerate}

\medskip
\noindent\textbf{Part III --- De Moivre's shortcuts.}
\begin{enumerate}[resume]
  \item The construction verdict: $\cos 72^\circ$ needs only
        square roots (question 9), so the pentagon is
        ruler-and-compass constructible; $\cos 20^\circ$
        satisfies an unbreakable cubic
        (\cref{pb:g12:trigo:1}), so the $60^\circ$ angle cannot
        be trisected. State the emerging criterion (square
        roots buildable, cube roots not), whose full theory ---
        Gauss, age 19, and the 17-gon --- lives in the
        university volumes.
  \item Solve $z^4 = -4$ by \omterm{def:g12:complex:argument}{exponential form} and recover
        \cref{exo:g12:complex:8}'s factorization
        $z^4 + 4 = (z^2 - 2z + 2)(z^2 + 2z + 2)$ by pairing
        \omterm{def:g12:complex:conjugate}{conjugate} \omterm{def:g11:quad:discriminant}{roots}.
  \item Expand $(\cos\theta + \iu\sin\theta)^3$ with the
        binomial theorem and de Moivre, and re-derive
        $\cos 3\theta = 4\cos^3\theta - 3\cos\theta$ in three
        lines --- compare with the addition-formula route of
        \cref{pb:g12:trigo:1}.
  \item Cross-check with \cref{exo:g12:complex:9}: for
        $\theta = \frac{2\pi}{5}$ and $n = 4$, what does the
        sum $1 + \cos\theta + \dots + \cos 4\theta$ equal, and
        why is that consistent with question 8?
  \item Compute the three cube roots of $8\iu$ in \omterm{def:g12:complex:def}{algebraic
        form}.
\end{enumerate}

\medskip
\noindent\textbf{Part IV --- Multiplication is geometry.}
\begin{enumerate}[resume]
  \item Describe completely the map $z \mapsto (1 + \iu)z$:
        by what angle does it rotate, by what factor does it
        scale, and what does it do to the unit square?
  \item Prove $\abs{z_1 z_2} = \abs{z_1}\,\abs{z_2}$ with
        \omterm{def:g12:complex:conjugate}{conjugates} (\cref{prop:g12:complex:modulus}), then
        compute $(1+\iu)^n$ for $n = 1, 2, 3, 4$ and describe
        the spiral the powers trace.
  \item The number $\mathrm{j} = \eu^{2\iu\pi/3}$: show
        $1 + \mathrm{j} + \mathrm{j}^2 = 0$, and verify on the triangle
        $\left(1, \mathrm{j}, \mathrm{j}^2\right)$ the classical criterion:
        $a + \mathrm{j} b + \mathrm{j}^2 c = 0$ holds for (one orientation
        of) an equilateral triangle $abc$.
  \item The fundamental theorem of algebra (admitted:
        d'Alembert--Gauss): every polynomial factors completely
        over $\C$. Deduce the real-world corollary --- every
        real polynomial factors into real linear and quadratic
        pieces --- and point to question 12 as an instance.
  \item Finale --- $\C$'s portrait, one sentence each: numbers
        become points; multiplication becomes
        rotation-and-scaling; the \omterm{def:g12:complex:rootsofunity}{roots of unity} become regular
        polygons; de Moivre turns powers into trigonometric
        identities; and algebraic closure guarantees that no
        \omterm{def:g10:algebra:equation}{equation} ever needs a bigger world. Coda: which two
        celebrated numbers met in question 10?
\end{enumerate}
\end{problem}
